\chapter{Affiene ruimten}\label{ch:b2:affine}

Vectorruimten hebben een bevoorrecht punt --- de oorsprong --- dat
de meetkunde niet wil. Een \emph{\omterm{def:b2:affine:def}{affiene ruimte}} is een vectorruimte
die haar oorsprong is vergeten: punten en vectoren worden
verschillende soorten, verbonden door translatie. Dit korte
hoofdstuk bouwt het woordenboek (punten, \omterm{def:b2:affine:barycenter}{barycentra}, \omterm{def:b2:affine:subspace}{affiene
deelruimten} en afbeeldingen), de \omterm{def:b2:affine:subspace}{affiene} kijk op convexiteit, en de
gereedschappen voor de klassering die in de meetkundige
hoofdstukken hierna worden gebruikt.

\section{Punten en vectoren}

\begin{definition}\label{def:b2:affine:def}
Een \emph{affiene ruimte}\index{affiene ruimte} gericht door een
reële vectorruimte $E$ is een niet-lege verzameling $\mathcal{E}$
met een afbeelding $(A, B) \mapsto \vect{AB} \in E$ die voldoet aan
\[
\vect{AB} + \vect{BC} = \vect{AC}
\quad \text{(Chasles)},
\qquad
\text{voor elke } A \text{ is } B \mapsto \vect{AB} \text{ een
bijectie } \mathcal{E} \to E .
\]
Men schrijft $B = A + u$ voor het unieke punt met $\vect{AB} = u$.
De \emph{dimensie} van $\mathcal{E}$ is $\dim E$. Elke vectorruimte
is een \omterm{def:b2:affine:subspace}{affiene} ruimte over zichzelf ($\vect{AB} = B - A$); elke
keuze van een oorsprong $O \in \mathcal{E}$ identificeert
$\mathcal{E}$ met $E$ via $M \mapsto \vect{OM}$.
\end{definition}

\begin{example}[Een affiene ruimte zonder natuurlijke
oorsprong]\label{ex:b2:affine:planeexample}
Het oplossingsvlak $\mathcal E = \{(x, y, z) \in \R^3 : x + y + z =
1\}$ is geen vectordeelruimte ($0 \notin \mathcal E$), maar wel een
\omterm{def:b2:affine:def}{affiene ruimte} gericht door $E = \{x + y + z = 0\}$: voor $A, B \in
\mathcal E$ belandt het verschil $\vect{AB} = B - A$ in $E$ (de
sommen heffen elkaar op), Chasles wordt van $\R^3$ geërfd, en $B
\mapsto \vect{AB}$ is bijectief op $E$. Geen enkel punt van
$\mathcal E$ is onderscheiden --- elke keuze van een ``oorsprong''
$O \in \mathcal E$ werkt even goed, en alle identificaties $M
\mapsto \vect{OM}$ verschillen slechts met translaties. Dit is de
typische toestand: oplossingsverzamelingen van inhomogene lineaire
problemen (lineaire stelsels, lineaire differentiaalvergelijkingen
in \cref{ch:b2:diffeq}) zijn \omterm{def:b2:affine:subspace}{affien}, nooit lineair, en de leuze
``particuliere oplossing plus kern'' is precies de uitspraak
$\mathcal F = A + F$ van de volgende definitie.
\end{example}

\begin{definition}[Barycentrum]\label{def:b2:affine:barycenter}
Zij $(A_i, \lambda_i)_{i \leq k}$ een familie gewogen punten met
$\sum\lambda_i \neq 0$. Het \emph{barycentrum}\index{barycentrum}
$G = \operatorname{bar}\bigl((A_i, \lambda_i)\bigr)$ is het unieke
punt met
\[
\sum_i \lambda_i\, \vect{GA_i} = 0
\qquad\text{gelijkwaardig}\qquad
\vect{OG} = \frac{1}{\sum\lambda_i}\sum_i \lambda_i\,\vect{OA_i}
\quad (\text{voor elke } O).
\]
Barycentra zijn \emph{associatief} (deelgroepen van punten mogen
worden vervangen door hun deelbarycentrum met het gesommeerde
gewicht) en invariant onder het herschalen van alle gewichten.
\end{definition}

\begin{proof}[Bewijs van het bestaan en van de formules]
Houd $O$ vast en schrijf $s = \sum_i\lambda_i \neq 0$. Volgens
Chasles is
\[
\sum_i\lambda_i\,\vect{GA_i} = 0
\iff \sum_i\lambda_i\bigl(\vect{GO} + \vect{OA_i}\bigr) = 0
\iff s\,\vect{OG} = \sum_i\lambda_i\,\vect{OA_i},
\]
wat $G = O + \frac1s\sum_i\lambda_i\vect{OA_i}$ eenduidig
vastlegt. \emph{Onafhankelijkheid van $O$:} voor een andere
oorsprong $O'$ is
\[
\frac1s\sum_i\lambda_i\,\vect{O'A_i}
= \frac1s\sum_i\lambda_i\bigl(\vect{O'O} + \vect{OA_i}\bigr)
= \vect{O'O} + \frac1s\sum_i\lambda_i\,\vect{OA_i}
= \vect{O'G} :
\]
hetzelfde punt $G$. \emph{Associativiteit:} splits de
indexverzameling als $I \sqcup J$ met $s_I = \sum_{i\in I}\lambda_i
\neq 0$, en zij $G_I$ het \omterm{def:b2:affine:barycenter}{barycentrum} van $(A_i, \lambda_i)_{i\in
I}$, zodat $\sum_{i\in I}\lambda_i\vect{OA_i} = s_I\,\vect{OG_I}$.
Dan is
\[
s\,\vect{OG}
= \sum_{i\in I}\lambda_i\vect{OA_i} +
\sum_{j\in J}\lambda_j\vect{OA_j}
= s_I\,\vect{OG_I} + \sum_{j\in J}\lambda_j\vect{OA_j} :
\]
$G$ is het \omterm{def:b2:affine:barycenter}{barycentrum} van $(G_I, s_I)$ samen met $(A_j,
\lambda_j)_{j\in J}$, zoals beweerd. \emph{Herschalen:} elke
$\lambda_i$ vervangen door $t\lambda_i$ ($t \neq 0$)
vermenigvuldigt $s$ en de gewogen som met $t$, en laat $\vect{OG}$
onveranderd.
\end{proof}

\begin{remark}[Klassieke valkuilen]
Twee vallen omringen de definitie. Ten eerste: sommeren de
gewichten tot \emph{nul}, dan is er geen \omterm{def:b2:affine:barycenter}{barycentrum}: de afbeelding
$O \mapsto \sum\lambda_i\vect{OA_i}$ is dan onafhankelijk van $O$
en definieert een \emph{vector}, geen punt --- bijvoorbeeld codeert
$(A, -1;\ B, 1)$ de vector $\vect{AB}$. Bijhouden welk van de twee
objecten een berekening voortbrengt, is de helft van de
barycentrische hygiëne. Ten tweede: gewichten hebben alleen
betekenis op een gemeenschappelijke factor ongelijk aan nul na;
formules als ``de coördinaten van $G$ zijn $\lambda_1, \dots,
\lambda_k$'' veronderstellen een normalisatie (gewoonlijk
$\sum\lambda_i = 1$), en vergeten te normaliseren is de
standaardbron van verkeerde verhoudingen op een figuur.
\end{remark}

\begin{definition}[Affiene deelruimten; affiene afbeeldingen]\label{def:b2:affine:subspace}
Een \emph{affiene deelruimte} is een verzameling $\mathcal{F} = A +
F = \{A + u : u \in F\}$ met $F$ een vectordeelruimte (haar
\emph{richting}); gelijkwaardig: een niet-lege verzameling die
stabiel is onder \omterm{def:b2:affine:barycenter}{barycentra}. De affiene deelruimten van $\R^n$ zijn
precies de oplossingsverzamelingen van lineaire stelsels $MX = B$
(bachelorjaar 1: particuliere oplossing plus kern). Een afbeelding
$f \colon \mathcal{E} \to \mathcal{E}'$ heet
\emph{affien}\index{affiene afbeelding} wanneer zij \omterm{def:b2:affine:barycenter}{barycentra}
bewaart --- gelijkwaardig wanneer
\[
f(A + u) = f(A) + \varphi(u)
\]
voor een (unieke) lineaire afbeelding $\varphi = \vec f$, het
\emph{lineaire deel}. Affiene afbeeldingen van $\R^n$: $X \mapsto
MX + C$. Samenstellingen zijn affien met samengestelde lineaire
delen; $f$ is bijectief dan en slechts dan als $\vec f$ dat is.
\end{definition}

\begin{proof}[Bewijs van de gelijkwaardigheid voor afbeeldingen]
Geldt $f(A + u) = f(A) + \varphi(u)$: voor een \omterm{def:b2:affine:barycenter}{barycentrum} $G$ van
$(A_i, \lambda_i)$ geeft alle punten vanuit $A$ uitwerken dat $f(G)
= f(A) + \varphi(\vect{AG})$ met $\varphi(\vect{AG}) =
\frac{\sum\lambda_i\varphi(\vect{AA_i})}{\sum\lambda_i}$: $f(G)$ is
het \omterm{def:b2:affine:barycenter}{barycentrum} van de beelden. Omgekeerd, houd $A$ vast en
definieer $\varphi(u) = \vect{f(A)\,f(A + u)}$.
\emph{Homogeniteit:} $A + tu = \operatorname{bar}\bigl(A, 1-t;\ A +
u, t\bigr)$ voor elke reële $t$, dus het bewaren van \omterm{def:b2:affine:barycenter}{barycentra}
(met willekeurige reële gewichten, zoals verondersteld) geeft
rechtstreeks $\varphi(tu) = t\,\varphi(u)$. \emph{Additiviteit:} $A
+ u + v = \operatorname{bar}\bigl(A + 2u, \tfrac12;\ A + 2v,
\tfrac12\bigr)$, dus $\varphi(u + v) = \tfrac12\varphi(2u) +
\tfrac12\varphi(2v) = \varphi(u) + \varphi(v)$, met de
homogeniteit. Bijgevolg is $\varphi$ lineair.
\end{proof}

\begin{remark}
Het bewijs gebruikte \omterm{def:b2:affine:barycenter}{barycentra} met \emph{willekeurige reële}
gewichten: de stap met de homogeniteit neemt $t$ buiten
$\intcc01$. Wordt van een afbeelding alleen verondersteld dat zij
\omterm{def:b2:affine:barycenter}{barycentra} met niet-negatieve gewichten bewaart --- gelijkwaardig:
middelpunten en segmenten --- dan komt de lineariteit van de
vectorafbeelding niet meer gratis: men krijgt slechts
$\Q$-lineariteit, en er is een hypothese van \omterm{def:b2:metric:continuity}{continuïteit} nodig om
te besluiten, precies als in \cref{exo:b2:affine:5}. ``Alle
\omterm{def:b2:affine:barycenter}{barycentra} bewaren'' onderscheiden van ``\omterm{def:b2:affine:convex}{convexe} combinaties
bewaren'' is een kleine maar echte subtiliteit van de \omterm{def:b2:affine:subspace}{affiene}
woordenschat.
\end{remark}

\begin{example}[Klassieke barycentrische meetkunde]\label{ex:b2:affine:median}
Het zwaartepunt van een driehoek $ABC$ is het \omterm{def:b2:affine:barycenter}{barycentrum} $G =
\operatorname{bar}(A,1; B,1; C,1)$. De associativiteit met het
middelpunt $A' = \operatorname{bar}(B, 1; C, 1)$ toont dat
\[
G = \operatorname{bar}(A, 1;\ A', 2) :
\]
$G$ ligt op de zwaartelijn $AA'$ op twee derde ervan --- en evenzo
voor de twee andere zwaartelijnen: de drie zwaartelijnen gaan door
één punt, in één regel barycentrische rekenkunde.
\end{example}

\begin{example}[De bimediaan van een
vierhoek]\label{ex:b2:affine:bimedians}
Zij $ABCD$ een willekeurige vierhoek (vlak of niet!) en beschouw
haar \emph{bimedianen}: de segmenten die de middelpunten van
overstaande zijden verbinden, $M_{AB}M_{CD}$ en $M_{BC}M_{DA}$.
Voer het \omterm{def:b2:affine:barycenter}{barycentrum} $G$ van $(A,1; B,1; C,1; D,1)$ in en groepeer
de gewichten op twee manieren:
\[
G = \operatorname{bar}\bigl(M_{AB}, 2;\ M_{CD}, 2\bigr)
= \operatorname{bar}\bigl(M_{BC}, 2;\ M_{DA}, 2\bigr) :
\]
$G$ is het middelpunt van \emph{beide} bimedianen --- de twee
bimedianen delen elkaar dus altijd middendoor, en de vierhoek van
de vier middelpunten is een parallellogram (haar diagonalen zijn de
bimedianen). Geen gevalsonderscheid, geen coördinaten, en het
argument blijft woordelijk overeind voor een scheve vierhoek in
$\R^3$, waar een bewijs op een tekening al kies zou zijn: de
associativiteit trekt zich niets van de dimensie aan.
\end{example}

\begin{example}[Een affiene afbeelding klasseren, van begin tot
eind]\label{ex:b2:affine:dilation}
Zij $f(x, y) = (2x - 1,\ 3y - 4)$ op $\R^2$. Haar lineaire deel is
$\varphi = \operatorname{diag}(2, 3)$, waarvan het spectrum $\{2,
3\}$ de waarde $1$ vermijdt: volgens het criterium voor vaste
punten dat hieronder wordt bewezen
(\cref{prop:b2:affine:fixedpoint}) heeft $f$ precies één vast punt,
gevonden door
\[
x = 2x - 1, \qquad y = 3y - 4
\qquad\Longrightarrow\qquad \Omega = (1, 2)
\]
op te lossen. Hercentreren in $\Omega$ (zet $x = 1 + u$, $y = 2 +
v$):
\[
f(1 + u,\ 2 + v) = (1 + 2u,\ 2 + 3v) :
\]
in het assenstelsel met oorsprong $\Omega$ \emph{is} $f$ haar
lineaire deel, een anisotrope uitrekking die vanuit het middelpunt
$(1, 2)$ horizontaal met $2$ en verticaal met $3$ vermenigvuldigt.
De algemene les: een \omterm{def:b2:affine:subspace}{affiene afbeelding} is ``lineaire afbeelding
plus plaatsgegevens'', en die plaatsgegevens klappen ineen tot één
goedgekozen oorsprong zodra $1$ geen \omterm{def:b2:reduction:eigen}{eigenwaarde} is. Omgekeerd
\emph{maakt} de oorsprong slecht verplaatsen de constante termen:
\omterm{def:b2:affine:subspace}{affiene} meetkunde is de kunst van het kiezen waar je $0$ legt.
\end{example}

\begin{remark}[Methode: samenloop en collineariteit via
barycentra]
\cref{ex:b2:affine:median} is een geval van een algemeen recept. Om
te bewijzen dat drie hoektransversalen van een driehoek door één
punt gaan, geef je één \emph{enkel} gewogen stelsel $(A, \alpha; B,
\beta; C, \gamma)$ en gebruik je de associativiteit op drie
manieren: $(B, C)$ groeperen toont dat het \omterm{def:b2:affine:barycenter}{barycentrum} op de
transversale uit $A$ ligt, $(C, A)$ groeperen dat het op die uit
$B$ ligt, en $(A, B)$ groeperen op de derde. Voor de zwaartelijnen
doet het stelsel $(A, 1; B, 1; C, 1)$ al het werk; voor
transversalen die de zijden in voorgeschreven verhoudingen snijden,
leest men de gewichten van die verhoudingen af. Om te bewijzen dat
drie punten \emph{op één rechte liggen}, schrijf je er één als
\omterm{def:b2:affine:barycenter}{barycentrum} van de twee andere (\cref{exo:b2:affine:2}), of gebruik
je het determinantcriterium van \cref{exo:b2:affine:11}. Beide
recepten vervangen meetkundig vernuft door boekhouding met
gewichten --- en daar dient de barycentrische rekenkunde precies
voor.
\end{remark}

\section{Convexiteit, affien bekeken}

\begin{definition}\label{def:b2:affine:convex}
Een deelverzameling $C$ van een \omterm{def:b2:affine:def}{affiene ruimte} heet
\emph{convex}\index{convexe verzameling} wanneer zij elk
\omterm{def:b2:affine:barycenter}{barycentrum} met \emph{niet-negatieve} gewichten van haar punten
bevat --- gelijkwaardig: elk segment $\intcc{A}{B} =
\{\operatorname{bar}(A, 1-t; B, t) : t \in \intcc{0}{1}\}$ tussen
haar punten. Het \emph{convexe omhulsel}\index{convexe omhulsel}
$\operatorname{conv}(S)$ is de verzameling van alle \omterm{def:b2:affine:barycenter}{barycentra} met
niet-negatieve gewichten van punten van $S$ --- de kleinste convexe
verzameling die $S$ bevat.
\end{definition}

\begin{example}[Epigrafen zijn convexe
verzamelingen]\label{ex:b2:affine:epigraph}
Het gebied $C = \{(x, y) : y \geq x^2\}$ boven de parabool is
\omterm{def:b2:affine:convex}{convex}: voor $(x_1, y_1), (x_2, y_2) \in C$ en $t \in \intcc01$
geeft de convexiteitsongelijkheid van de kwadraatfunctie
\[
\bigl((1-t)x_1 + tx_2\bigr)^2 \leq (1-t)x_1^2 + tx_2^2
\leq (1-t)y_1 + ty_2 ,
\]
zodat het \omterm{def:b2:affine:barycenter}{barycentrum} boven de parabool blijft. De berekening is
algemeen: $\{y \geq f(x)\}$ is \omterm{def:b2:affine:convex}{convex} precies wanneer $f$ een
convexe functie is --- \omterm{def:b2:affine:convex}{convexe} \emph{verzamelingen} en \omterm{def:b2:affine:convex}{convexe}
\emph{functies} (\cref{ch:b2:realfun}) zijn twee gezichten van één
begrip, met de epigrafen als woordenboek. Dit is de meetkundige
reden waarom convexe functies steunrechten hebben, het feit dat in
\cref{ch:b2:randomvar} de ongelijkheid van Jensen zal bewijzen.
\end{example}

\begin{example}[Overbodige voortbrengers]\label{ex:b2:affine:squarehull}
Zij $S = \{(0,0), (2,0), (2,2), (0,2), (1,1)\}$. Het vijfde punt is
het \omterm{def:b2:affine:barycenter}{barycentrum}
\[
(1,1) = \operatorname{bar}\bigl((0,0), \tfrac12;\ (2,2),
\tfrac12\bigr),
\]
en ligt dus al in het omhulsel van de vier andere:
$\operatorname{conv}(S)$ is het vierkant met de vier hoeken als
hoekpunten. In het algemeen mag een punt van $S$ dat een
\omterm{def:b2:affine:barycenter}{barycentrum} met niet-negatieve gewichten van de \emph{overige}
punten van $S$ is, worden geschrapt zonder het omhulsel te
veranderen; de punten die nooit geschrapt kunnen worden (hier de
vier hoeken) zijn de \emph{extreme punten} van het omhulsel. Ze
bepalen is een zuivere berekening met \omterm{def:b2:affine:barycenter}{barycentra}: $(2,0)$
bijvoorbeeld laat zich niet als $\operatorname{bar}$ van de overige
punten met niet-negatieve gewichten schrijven, omdat de eerste
coördinaat al het gewicht op punten met $x = 2$ zou leggen,
waarna de tweede coördinaat faalt. Vragen over convexiteit
herleiden telkens opnieuw tot het oplossen van kleine gewogen
stelsels.
\end{example}

\begin{theorem}[Carathéodory]\label{thm:b2:affine:caratheodory}
In een \omterm{def:b2:affine:def}{affiene ruimte} van dimensie $n$ is elk punt van
$\operatorname{conv}(S)$ een \omterm{def:b2:affine:barycenter}{barycentrum} van hoogstens $n + 1$
punten van $S$.\index{stelling van Carathéodory}
\end{theorem}

\begin{proof}
Zij $G = \operatorname{bar}(A_0, \lambda_0; \dots; A_k,
\lambda_k)$ met $\lambda_i > 0$, $\sum\lambda_i = 1$ en $k + 1 > n
+ 1$ punten. De $k$ vectoren $\vect{A_0A_i}$ ($i \geq 1$) zijn
afhankelijk ($k > n$): $\sum_{i\geq1}\mu_i \vect{A_0A_i} = 0$ op
niet-triviale wijze; zetten we $\mu_0 = -\sum_{i\geq1}\mu_i$, dan
krijgen we gewichten $(\mu_i)$ met $\sum\mu_i = 0$ en $\sum
\mu_i\,\vect{OA_i} = 0$ (voor elke $O$), niet alle nul. Dan
sommeren de gewichten $\lambda_i - t\mu_i$ voor elke reële $t$ nog
altijd tot $1$ en geldt, omdat $\sum_i\mu_i\vect{OA_i} = 0$,
\[
\sum_i(\lambda_i - t\mu_i)\,\vect{OA_i}
= \sum_i\lambda_i\,\vect{OA_i} :
\]
zij brengen \emph{hetzelfde} punt $G$ voort. Laat $t$ nu vanaf $0$
lopen: een zekere $\mu_i$ is positief (zij sommeren tot nul en zijn
niet alle nul), dus is
\[
t^* = \min\Bigl\{\frac{\lambda_i}{\mu_i} : \mu_i > 0\Bigr\}
\]
welgedefinieerd en positief. In $t = t^*$: voor de indices met
$\mu_i > 0$ is $\lambda_i - t^*\mu_i \geq 0$ wegens de
minimaliteit, met gelijkheid in een minimaliserende index; voor de
indices met $\mu_i \leq 0$ is $\lambda_i - t^*\mu_i \geq \lambda_i
> 0$. Alle gewichten blijven niet-negatief en minstens één is
gestorven: $G$ is herschreven als \omterm{def:b2:affine:barycenter}{barycentrum} van minder punten.
Herhaal zolang er meer dan $n + 1$ punten overblijven.
\end{proof}

\begin{example}
In het vlak ($n = 2$): elk punt van het \omterm{def:b2:affine:convex}{convexe omhulsel} van een
eindige verzameling ligt in een driehoek met hoekpunten in die
verzameling --- de meetkundige inhoud van Carathéodory, evenzeer
gebruikt in de optimalisatie als in de kansrekening (mengsels).
\end{example}

\begin{example}[Het algoritme van Carathéodory
uitvoeren]\label{ex:b2:affine:caratheodoryrun}
Schrijf het middelpunt van het vierkant uit
\cref{ex:b2:affine:squarehull} met zijn vier hoeken $A_1 = (0,0)$,
$A_2 = (2,0)$, $A_3 = (2,2)$, $A_4 = (0,2)$:
\[
(1,1) = \operatorname{bar}\bigl(A_1, \tfrac14;\ A_2,
\tfrac14;\ A_3, \tfrac14;\ A_4, \tfrac14\bigr),
\]
vier punten in dimensie $2$ --- één te veel. Het recept van het
bewijs vraagt gewichten $(\mu_i)$ met $\sum\mu_i = 0$ en
$\sum\mu_i\vect{OA_i} = 0$: hier voldoet $\mu = (1, -1, 1, -1)$ (de
twee diagonalen delen hun middelpunt). De schuif $\lambda_i \mapsto
\lambda_i - t\mu_i$ houdt het \omterm{def:b2:affine:barycenter}{barycentrum} voor elke $t$ vast; de
uiterste toelaatbare waarde $t = \frac14$ maakt de gewichten $(0,
\tfrac12, 0, \tfrac12)$ en doodt $A_1$ en $A_3$ tegelijk:
\[
(1,1) = \operatorname{bar}\bigl(A_2, \tfrac12;\ A_4,
\tfrac12\bigr),
\]
een voorstelling met twee punten --- zelfs beter dan de drie die de
stelling waarborgt, omdat het middelpunt toevallig op een segment
tussen voortbrengers ligt. Het algoritme is volstrekt mechanisch:
zoek een afhankelijkheid, schuif tot een gewicht sterft, herhaal.
\end{example}

\section{Gereedschappen voor de affiene klassering}

\begin{proposition}[Vaste punten van affiene afbeeldingen]\label{prop:b2:affine:fixedpoint}
Zij $f$ een \omterm{def:b2:affine:subspace}{affien} endomorfisme van een eindigdimensionale \omterm{def:b2:affine:def}{affiene
ruimte} met lineair deel $\varphi$. Geldt $1 \notin
\operatorname{Sp}(\varphi)$, dan heeft $f$ precies één vast punt
$\Omega$, en in de vectorialisatie in $\Omega$ \emph{is} $f$ haar
lineaire deel. (Translaties, met $\varphi = \mathrm{id}$ en zonder
vast punt, zijn de basisobstructie.)
\end{proposition}

\begin{proof}
Houd $O$ vast en schrijf $f(O + x) = f(O) + \varphi(x)$. Het punt
$O + x$ is vast dan en slechts dan als $O + x = f(O) + \varphi(x)$,
dat wil zeggen
\[
(\mathrm{id} - \varphi)(x) = \vect{O f(O)} .
\]
In eindige dimensie is $\mathrm{id} - \varphi$ inverteerbaar dan en
slechts dan als $0$ geen \omterm{def:b2:reduction:eigen}{eigenwaarde} van $\mathrm{id} - \varphi$
is, dat wil zeggen dan en slechts dan als $1 \notin
\operatorname{Sp}\varphi$ --- en in dat geval heeft de getoonde
vergelijking precies één oplossing $x^*$, wat het unieke vaste punt
$\Omega = O + x^*$ geeft. Hercentreren: voor elke vector $u$ is
\[
f(\Omega + u) = f(\Omega) + \varphi(u) = \Omega + \varphi(u),
\]
zodat de afbeelding in het assenstelsel met oorsprong $\Omega$
luidt als $u \mapsto \varphi(u)$: zuiver lineair. Geldt $1 \in
\operatorname{Sp}\varphi$, dan bestaat er ofwel geen vast punt (de
getoonde vergelijking kan onoplosbaar zijn, zoals voor een
translatie), ofwel een hele \omterm{def:b2:affine:subspace}{affiene deelruimte} vaste punten (tel
bij een oplossing een willekeurige \omterm{def:b2:reduction:eigen}{eigenvector} bij \omterm{def:b2:reduction:eigen}{eigenwaarde} $1$
op): de eenduidigheid is precies de spectrale voorwaarde.
\end{proof}

\begin{example}[Isometrieën van het vlak, voltooid]\label{ex:b2:affine:isometries}
Een \omterm{def:b2:affine:subspace}{affiene} isometrie van het euclidische vlak heeft haar lineaire
deel in $O(2)$: een rotatie $R_\theta$ of een spiegeling (volume
van bachelorjaar 1). Is $\theta \neq 0$, dan is $1 \notin
\operatorname{Sp} R_\theta$, dus is de afbeelding een
\emph{rotatie om een uniek middelpunt}
(\cref{prop:b2:affine:fixedpoint}). Is het lineaire deel een
spiegeling, dan is zij ofwel een spiegeling in een as (er bestaan
vaste punten) ofwel een \emph{schuifspiegeling} (een spiegeling
samengesteld met een translatie langs de as, zonder vast punt).
Samen met de translaties is dit de volledige klassering van de
isometrieën van het vlak.
\end{example}

\begin{remark}[De isometrieën van het vlak, in één oogopslag]
De gevallen verzameld: de identiteit; translaties ($\vec f =
\mathrm{id}$, zonder vast punt behalve in het triviale geval);
rotaties (lineair deel $R_\theta$ met $\theta \neq 0$: één
middelpunt); spiegelingen (lineair deel een spiegeling, een rechte
vaste punten); schuifspiegelingen (hetzelfde lineaire deel, geen
vast punt). Vier families plus de identiteit, elk herkend aan
slechts twee gegevens: het lineaire deel en de verzameling vaste
punten --- het patroon van \cref{prop:b2:affine:fixedpoint},
uitputtend gemaakt.
\end{remark}

\begin{example}[Een schuifspiegeling, op heterdaad
betrapt]\label{ex:b2:affine:glide}
Zij $f(x, y) = (y + 1,\ x + 1)$. Het lineaire deel $(x, y) \mapsto
(y, x)$ is de spiegeling in de diagonaal $y = x$, dus is $1 \in
\operatorname{Sp}\vec f$ en zwijgt
\cref{prop:b2:affine:fixedpoint}. Vaste punten zouden tegelijk $x =
y + 1$ en $y = x + 1$ vereisen: onmogelijk --- er zijn er geen, dus
$f$ is geen spiegeling. Kwadrateren beslist de klassering:
\[
f\bigl(f(x, y)\bigr) = f(y + 1,\ x + 1) = (x + 2,\ y + 2),
\]
de translatie over $(2, 2)$: $f$ is de \emph{schuifspiegeling} met
als as de rechte $y = x$ (passend verschoven: het middelpunt van
$M$ en $f(M)$ ligt altijd op $y = x + {}$constante, hier $y = x$,
zoals men op $M = (0, 0) \mapsto (1,1)$ nagaat) en met
schuifvector $(1, 1)$, de helft van $f \circ f$. Vergelijk met
\cref{exo:b2:affine:6}, waar hetzelfde lineaire deel maar een
andere constante een echte spiegeling opleverde: is de \omterm{def:b2:reduction:eigen}{eigenwaarde}
$1$ aanwezig, dan beslist de constante term alles.
\end{example}

\begin{example}[Affiene recursies zijn affiene
dynamica]\label{ex:b2:affine:recursion}
De klassieke recursie $u_{n+1} = au_n + b$ ($a \neq 1$) itereert de
\omterm{def:b2:affine:subspace}{affiene afbeelding} $f(x) = ax + b$ van de rechte, waarvan het
lineaire deel $a$ de \omterm{def:b2:reduction:eigen}{eigenwaarde} $1$ vermijdt: er is een uniek vast
punt $\omega = \frac{b}{1-a}$, en daar hercentreren (het
eendimensionale geval van de propositie hierboven) maakt van $f$
een vermenigvuldiging met $a$:
\[
u_{n+1} - \omega = a\,(u_n - \omega)
\qquad\Longrightarrow\qquad
u_n = \omega + a^n(u_0 - \omega) .
\]
Voor $u_{n+1} = \frac{u_n}2 + 3$: $\omega = 6$ en $u_n = 6 + (u_0 -
6)2^{-n} \to 6$. Het recept dat in \cref{ch:b2:series} voor zulke
recursies wordt geleerd --- ``trek het vaste punt af'' --- is
precies de vectorialisatie van een \omterm{def:b2:affine:subspace}{affiene afbeelding} in haar vaste
punt; de convergentie voor $\abs a < 1$ is het
contractieverschijnsel dat \cref{ch:b2:metric} tot de
vastepuntsstelling van Banach maakte. Eén idee, drie hoofdstukken.
\end{example}

\begin{example}[Het middelpunt van een rotatie vinden]\label{ex:b2:affine:rotationcenter}
Zij $f(x, y) = (-y + 2,\ x)$. Het lineaire deel is $\varphi(x, y) =
(-y, x)$: de rotatie over de hoek $\frac\pi2$, waarvan het spectrum
$\{\iu, -\iu\}$ de waarde $1$ vermijdt. Volgens
\cref{prop:b2:affine:fixedpoint} is er precies één vast punt: $x =
-y + 2$ en $y = x$ geven $x = 1$, $y = 1$, dus $\Omega = (1, 1)$,
en $f$ \emph{is} de rotatie met middelpunt $(1, 1)$ over de hoek
$\frac\pi2$. De algemene les: is $1 \notin \operatorname{Sp}\vec
f$, dan kost het klasseren van $f$ één lineair stelsel --- de
meetkunde zit \omterm{def:b2:metric:complete}{volledig} in het lineaire deel, de rekenkunde \omterm{def:b2:metric:complete}{volledig}
in het lokaliseren van het middelpunt.
\end{example}

\begin{remark}[Waar de affiene taal hierna wordt gebruikt]
\omterm{def:b2:affine:barycenter}{Barycentra} en \omterm{def:b2:affine:subspace}{affiene afbeeldingen} zijn de grammatica van de
meetkundige hoofdstukken hierna: raaklijnen en raakvlakken zijn
\omterm{def:b2:affine:subspace}{affiene} objecten (\cref{ch:b2:curves,ch:b2:surfaces}), een \omterm{def:b2:affine:subspace}{affiene}
verandering van veranderlijke vermenigvuldigt oppervlakten en
volumes met $\abs{\det \vec f}$ (\cref{ch:b2:multint}), en de
verwachtingswaarde is een \omterm{def:b2:affine:barycenter}{barycentrum} met gewichten die door een
kansverdeling worden gegeven, en daarom beheerst de convexiteit de
ongelijkheid van Jensen (\cref{ch:b2:randomvar}). In het volume van
bachelorjaar 3 draagt dezelfde woordenschat van de convexiteit de
studie van de $L^p$-normen en van de integraalongelijkheden.
\end{remark}

\begin{remark}[Vooruitblik binnen dit volume]
Twee draden verlaten dit hoofdstuk. De \emph{\omterm{def:b2:affine:subspace}{affiene}} draad:
raaklijnen (\cref{ch:b2:curves}) en raakvlakken
(\cref{ch:b2:surfaces}) zijn \omterm{def:b2:affine:subspace}{affiene deelruimten} die aan
niet-lineaire objecten hangen, en de klassering van de kwadrieken
in het hoofdstuk over oppervlakken draait op de
middelpuntsvergelijking $A\Omega = -b$ van dit hoofdstuk. De
\emph{\omterm{def:b2:affine:convex}{convexe}} draad is langer: de convexiteit van halfvlakken en
schijven drijft de theorie van Helly in de weekendopgave aan; de
convexiteit van functies geeft de ongelijkheid van Jensen
(\cref{ch:b2:randomvar}); en de laatste stelling van het boek ---
het uitstervingscriterium voor vertakkingsprocessen
(\cref{ch:b2:genfun}) --- wordt beslist door de ligging van een
convexe kromme ten opzichte van de diagonaal, een beeld dat evenzeer
bij dit hoofdstuk hoort als bij de kansrekening. Ook \omterm{def:b2:affine:barycenter}{barycentra}
keren daar terug: een verwachtingswaarde is een \omterm{def:b2:affine:barycenter}{barycentrum} met
kansgewichten.
\end{remark}

\section{Oefeningen}

\begin{exercise}[$\star$]\label{exo:b2:affine:1}
Zijn de volgende verzamelingen \omterm{def:b2:affine:subspace}{affiene deelruimten} van $\R^3$? Geef
richtingen en dimensies. $\{x + y + z = 1\}$; $\;\{x + y + z = 1,\
x - z = 3\}$; $\;\{x^2 + y^2 = 1\}$; de oplossingsverzameling van
$MX = B$ voor een gegeven verenigbaar stelsel.
\end{exercise}

\begin{exercise}[$\star$]\label{exo:b2:affine:2}
Bewijs dat drie verschillende punten $A, B, C$ van een \omterm{def:b2:affine:def}{affiene
ruimte} op één rechte liggen dan en slechts dan als $C$ een
\omterm{def:b2:affine:barycenter}{barycentrum} van $A$ en $B$ is, en dan en slechts dan als de
vectoren $\vect{AB}, \vect{AC}$ afhankelijk zijn. Leid daaruit
boekhouding met gewichten in de stijl van Menelaos af: is $C =
\operatorname{bar}(A, 1 - t; B, t)$, lokaliseer $C$ dan voor $t =
\frac12$, $t = 2$ en $t = -1$.
\end{exercise}

\begin{exercise}[$\star$]\label{exo:b2:affine:3}
Zij $f$ de \omterm{def:b2:affine:subspace}{affiene afbeelding} van $\R^2$ gegeven door $f(X) = MX +
C$ met $M = \frac12\begin{pmatrix} 1 & 1\\ 1 & 1\end{pmatrix}$ en
$C = (1, 0)^{\mathsf T}$. Bepaal het beeld van $f$, haar vaste
punten (als die er zijn) en $f \circ f$.
\end{exercise}

\begin{exercise}[$\star\star$]\label{exo:b2:affine:4}
(Associativiteit in actie) Verdeel in een driehoek $ABC$ de zijden
$BC$, $CA$, $AB$ met $I, J, K$ in de verhoudingen $\vect{BI} =
\frac13\vect{BC}$, $\vect{CJ} = \frac13\vect{CA}$, $\vect{AK} =
\frac13\vect{AB}$. Druk $I, J, K$ uit als \omterm{def:b2:affine:barycenter}{barycentra} en bereken het
\omterm{def:b2:affine:barycenter}{barycentrum} van $(I,1;J,1;K,1)$: wat vind je, en waarom was dat te
voorzien?
\end{exercise}

\begin{exercise}[$\star\star$]\label{exo:b2:affine:5}
Bewijs dat een afbeelding $f \colon \R^n \to \R^n$ die
\emph{middelpunten} bewaart ($f\bigl(\frac{A+B}{2}\bigr) =
\frac{f(A) + f(B)}{2}$) en \emph{\omterm{def:b2:metric:continuity}{continu}} is, \omterm{def:b2:affine:subspace}{affien} is.
\emph{(Toon aan dat de vectorafbeelding $u \mapsto f(O + u) - f(O)$
additief is via middelpunten, dan $\Q$-homogeen, dan
$\R$-homogeen wegens de \omterm{def:b2:metric:continuity}{continuïteit} --- dezelfde
dichtheidsstrategie als voor de functionaalvergelijking van Cauchy
in het volume van bachelorjaar 1; leid de nodige stappen hier
opnieuw af.)}
\end{exercise}

\begin{exercise}[$\star\star$]\label{exo:b2:affine:6}
Klasseer de \omterm{def:b2:affine:subspace}{affiene afbeelding} $f(x, y) = (y + 1,\; x - 1)$ van het
euclidische vlak: lineair deel, vaste punten, meetkundige aard
(spiegeling? schuifspiegeling?). Bereken $f \circ f$ en besluit.
\end{exercise}

\begin{exercise}[$\star\star\star$]\label{exo:b2:affine:7}
(Radon) Zijn $A_1, \dots, A_{n+2}$ punten van een \omterm{def:b2:affine:def}{affiene ruimte}
van dimensie $n$. Bewijs dat zij in twee disjuncte groepen kunnen
worden gesplitst waarvan de \omterm{def:b2:affine:convex}{convexe omhulsels} elkaar snijden.
\emph{(Zoek, als in het bewijs van Carathéodory, gewichten $\mu_i$,
niet alle nul, met $\sum\mu_i = 0$ en $\sum\mu_i\vect{OA_i} = 0$;
scheid de positieve van de negatieve gewichten en normaliseer beide
kanten.)}
\end{exercise}

\begin{exercise}[$\star\star\star$]\label{exo:b2:affine:8}
Zij $f$ een \omterm{def:b2:affine:subspace}{affien} endomorfisme van $\R^n$ met $f \circ f = f$.
Bewijs dat $f$ de \omterm{def:b2:affine:subspace}{affiene} projectie is op de \omterm{def:b2:affine:subspace}{affiene deelruimte}
$\operatorname{Fix}(f) = \operatorname{im} f$ langs de richting
$\ker\vec f$, en omgekeerd dat al zulke projecties idempotent zijn.
\emph{(Toon eerst aan dat $\operatorname{im} f$ uit vaste punten
bestaat.)}
\end{exercise}

\begin{exercise}[$\star$]\label{exo:b2:affine:9}
Zij $G = \operatorname{bar}(A, 1;\ B, 2;\ C, 3)$ in een driehoek
$ABC$. Toon met de associativiteit aan dat de rechte $AG$ de zijde
$BC$ snijdt in $M = \operatorname{bar}(B, 2;\ C, 3)$, en lokaliseer
$G$ op het segment $\intcc AM$; lokaliseer op dezelfde manier het
snijpunt van $BG$ met $CA$.
\end{exercise}

\begin{exercise}[$\star\star$]\label{exo:b2:affine:10}
Voor $\lambda \neq 0$ is de \emph{homothetie} $h_{\Omega, \lambda}$
de \omterm{def:b2:affine:subspace}{affiene afbeelding} die $\Omega$ vasthoudt met lineair deel
$\lambda\,\mathrm{id}$. Bewijs dat de samenstelling $h_{\Omega',
\mu} \circ h_{\Omega, \lambda}$ een homothetie met verhouding
$\lambda\mu$ is wanneer $\lambda\mu \neq 1$, en een translatie
wanneer $\lambda\mu = 1$; bereken in het geval $\lambda = \mu = -1$
(twee puntspiegelingen) de translatievector.
\end{exercise}

\begin{exercise}[$\star\star$]\label{exo:b2:affine:11}
(Menelaos) Zij in een driehoek $ABC$ $A' \in (BC)$, $B' \in (CA)$,
$C' \in (AB)$, alle verschillend van de hoekpunten, en definieer
$\alpha, \beta, \gamma$ door $\vect{A'B} = \alpha\,\vect{A'C}$,
$\vect{B'C} = \beta\,\vect{B'A}$, $\vect{C'A} =
\gamma\,\vect{C'B}$. Bewijs dat $A', B', C'$ op één rechte liggen
dan en slechts dan als $\alpha\beta\gamma = 1$. \emph{(Schrijf elk
punt als \omterm{def:b2:affine:barycenter}{barycentrum} van twee hoekpunten; toon aan dat drie punten
op één rechte liggen dan en slechts dan als hun rijen
barycentrische coördinaten ten opzichte van $(A, B, C)$ een
singuliere $3 \times 3$-matrix vormen.)}
\end{exercise}

\begin{exercise}[$\star\star\star$]\label{exo:b2:affine:12}
Bewijs dat het \omterm{def:b2:affine:convex}{convexe omhulsel} van een \omterm{def:b2:metric:compact}{compacte} deelverzameling
$K$ van $\R^n$ \omterm{def:b2:metric:compact}{compact} is. \emph{(Volgens
\cref{thm:b2:affine:caratheodory} is $\operatorname{conv}(K)$ het
beeld van een \omterm{def:b2:metric:compact}{compacte} verzameling onder een \omterm{def:b2:metric:continuity}{continue}
afbeelding.)} Toon met een voorbeeld in $\R^2$ aan dat het \omterm{def:b2:affine:convex}{convexe
omhulsel} van een \emph{gesloten} verzameling niet gesloten hoeft te
zijn.
\end{exercise}

\section{Probleem: van Radon naar Helly, centrumpunten en de
stelling van Jung}

\begin{omfigure}
\begin{tikzpicture}[scale=0.9]
  \draw[omDef, thick] (0,0) -- (3,0) -- (0.9,2.6) -- cycle;
  \fill (0,0) circle (1.6pt) node[below left] {\small $A_1$};
  \fill (3,0) circle (1.6pt) node[below right] {\small $A_2$};
  \fill (0.9,2.6) circle (1.6pt) node[above] {\small $A_3$};
  \fill[omThm] (1.3,0.9) circle (1.8pt)
    node[below, omThm] {\small $A_4$};
\end{tikzpicture}
\qquad
\begin{tikzpicture}[scale=0.9]
  \draw[omDef, thick] (0,0) -- (3,0.3) -- (3.3,2.4) --
    (0.4,2.2) -- cycle;
  \draw[omThm, thick] (0,0) -- (3.3,2.4);
  \draw[omThm, thick] (3,0.3) -- (0.4,2.2);
  \fill (0,0) circle (1.6pt) node[below left] {\small $A_1$};
  \fill (3,0.3) circle (1.6pt) node[below right] {\small $A_2$};
  \fill (3.3,2.4) circle (1.6pt) node[above right] {\small $A_3$};
  \fill (0.4,2.2) circle (1.6pt) node[above left] {\small $A_4$};
  \fill[omProp] (1.75,1.27) circle (1.8pt);
\end{tikzpicture}

{\small De twee types van Radon voor vier punten van het vlak in
algemene ligging: één punt binnen de driehoek van de andere
(partitie $\{A_4\} \mid \{A_1, A_2, A_3\}$), of \omterm{def:b2:affine:convex}{convexe} ligging,
waarbij het punt van Radon (oranje) het snijpunt van de twee
diagonalen is.}
\end{omfigure}

\begin{problem}[{Weekendopgave --- de stelling van Helly en twee van
haar dividenden}]\label{pb:b2:affine:1}
Het lemma van Radon (\cref{exo:b2:affine:7}) zegt dat $n + 2$
punten van een $n$-dimensionale \omterm{def:b2:affine:def}{affiene ruimte} altijd in twee
groepen splitsen waarvan de \omterm{def:b2:affine:convex}{convexe omhulsels} elkaar snijden. Deze
opgave maakt van dat ene feit uit de lineaire algebra een ketting
van stellingen uit de combinatorische meetkunde: de
doorsnedestelling van Helly, de stelling van het centrumpunt (een
mediaan in twee dimensies) en de overdekkingsstelling van Jung.
Overal is het vlak $\R^2$ met zijn gewone euclidische structuur, en
is $\det$ de \omterm{def:b2:linalg:det}{determinant} in de canonieke basis.\index{stelling van
Helly}\index{lemma van Radon}\index{stelling van
Jung}\index{centrumpunt}

\medskip
\noindent\textbf{Deel I --- Barycentrische coördinaten.} Punten
$A_0, \dots, A_k$ heten \emph{\omterm{def:b2:affine:subspace}{affien} onafhankelijk} wanneer de
vectoren $\vect{A_0A_1}, \dots, \vect{A_0A_k}$ lineair
onafhankelijk zijn.
\begin{enumerate}
  \item Toon aan dat de \omterm{def:b2:affine:subspace}{affiene} onafhankelijkheid niet van de
        keuze van het basispunt $A_0$ afhangt, en dat zij
        gelijkwaardig is met: telkens wanneer twee families
        gewichten, elk sommerend tot $1$, hetzelfde \omterm{def:b2:affine:barycenter}{barycentrum}
        van $(A_0, \dots, A_k)$ definiëren, vallen de gewichten
        samen.
  \item Zij $(A, B, C)$ \omterm{def:b2:affine:subspace}{affien} onafhankelijk in het vlak. Toon aan
        dat elk punt $M$ een uniek drietal $(\alpha, \beta,
        \gamma)$ toelaat met $\alpha + \beta + \gamma = 1$ en $M =
        \operatorname{bar}(A, \alpha; B, \beta; C, \gamma)$ ---
        zijn \emph{barycentrische coördinaten}.
  \item Bewijs de determinantformules
        \[
        \alpha = \frac{\det(\vect{MB},
        \vect{MC})}{\det(\vect{AB}, \vect{AC})},
        \qquad
        \beta = \frac{\det(\vect{MC},
        \vect{MA})}{\det(\vect{AB}, \vect{AC})},
        \qquad
        \gamma = \frac{\det(\vect{MA},
        \vect{MB})}{\det(\vect{AB}, \vect{AC})} :
        \]
        barycentrische coördinaten zijn verhoudingen van
        georiënteerde oppervlakten.
  \item De rechten $BC$, $CA$, $AB$ zijn de coördinaatrechten
        $\{\alpha = 0\}$, $\{\beta = 0\}$, $\{\gamma = 0\}$. Toon
        aan dat $M$ in de gesloten driehoek $\operatorname{conv}\{A,
        B, C\}$ ligt dan en slechts dan als $\alpha, \beta, \gamma
        \geq 0$, en dat de drie rechten het vlak in precies zeven
        gebieden snijden, geklasseerd door de tekens van $(\alpha,
        \beta, \gamma)$ (waarbij het tekenpatroon $(-, -, -)$
        onmogelijk is).
  \item Zij $u \colon \R^2 \to \R$ een \omterm{def:b2:affine:subspace}{affiene afbeelding} (een
        \emph{\omterm{def:b2:affine:subspace}{affiene} vorm}). Toon aan dat $u(M) = \alpha\,u(A) +
        \beta\,u(B) + \gamma\,u(C)$, dat de niveauverzamelingen
        van een niet-constante \omterm{def:b2:affine:subspace}{affiene} vorm rechten zijn, dat elke
        rechte zo ontstaat, en dat de gesloten halfvlakken $\{u
        \geq c\}$ \omterm{def:b2:affine:convex}{convex} zijn.
\end{enumerate}

\medskip
\noindent\textbf{Deel II --- Partities van Radon, verfijnd.} Een
familie van $n + 2$ punten van $\R^n$ ligt in \emph{algemene
ligging} wanneer elke $n + 1$ ervan \omterm{def:b2:affine:subspace}{affien} onafhankelijk zijn. Een
\emph{\omterm{def:b2:affine:subspace}{affiene} afhankelijkheid} van $(A_1, \dots, A_{n+2})$ is een
familie $(\mu_i)$ met $\sum_i \mu_i = 0$ en $\sum_i
\mu_i\,\vect{OA_i} = 0$ voor één (en dus voor elke) oorsprong $O$.
\begin{enumerate}[resume]
  \item Bereken een \omterm{def:b2:affine:subspace}{affiene} afhankelijkheid ongelijk aan nul van de
        vier punten $A_1 = (0,0)$, $A_2 = (3,0)$, $A_3 = (0,3)$,
        $A_4 = (1,1)$; geef de partitie van Radon en het punt van
        Radon.
  \item Toon aan dat voor punten in algemene ligging de
        vectorruimte van de \omterm{def:b2:affine:subspace}{affiene} afhankelijkheden dimensie
        precies $1$ heeft, en dat een afhankelijkheid ongelijk aan
        nul \emph{geen} verdwijnende coëfficiënt heeft.
  \item Leid af dat de partitie van Radon van $n + 2$ punten in
        algemene ligging uniek is (op het verwisselen van de twee
        blokken na), waarbij elk blok de verzameling indices is
        waar $\mu_i$ één vast teken heeft.
  \item Toon voor vier punten van het vlak in algemene ligging de
        tweedeling aan: ofwel heeft de partitie het type $(1, 3)$
        --- één punt inwendig aan de driehoek van de drie andere
        --- ofwel het type $(2, 2)$: de vier punten liggen in
        \omterm{def:b2:affine:convex}{convexe} ligging en de segmenten die de twee paren
        verbinden (de diagonalen) snijden elkaar, in het punt van
        Radon.
  \item Voer vraag 6 uit voor het eenheidsvierkant $(0,0)$,
        $(1,0)$, $(1,1)$, $(0,1)$: afhankelijkheid, partitie, punt
        van Radon.
\end{enumerate}

\medskip
\noindent\textbf{Deel III --- De stelling van Helly in het vlak.}
\begin{enumerate}[resume]
  \item Zijn $C_1, C_2, C_3, C_4$ \omterm{def:b2:affine:convex}{convexe} deelverzamelingen van
        $\R^2$ waarvan elke drie een gemeenschappelijk punt
        hebben. Kies $x_i \in \bigcap_{j \neq i} C_j$ en pas het
        lemma van Radon toe op $x_1, \dots, x_4$: toon aan dat het
        punt van Radon tot alle vier de verzamelingen behoort.
        \emph{(Voor elke $k$ bestaat het blok dat $x_k$ niet
        bevat, uit punten van $C_k$.)}
  \item (Helly) Zijn $C_1, \dots, C_m$ ($m \geq 3$) \omterm{def:b2:affine:convex}{convexe}
        deelverzamelingen van $\R^2$ waarvan elke drie elkaar
        snijden. Bewijs dat $\bigcap_{i=1}^m C_i \neq \emptyset$,
        met inductie naar $m$: vervang $C_{m-1}$ en $C_m$ door
        $C_{m-1} \cap C_m$ en ga de hypothese voor de nieuwe
        familie na met vraag 11.
  \item Drie tegenvoorbeelden, één per hypothese: (a) de drie
        gesloten zijden van een driehoek snijden elkaar
        paarsgewijs maar hebben geen gemeenschappelijk punt ($3$
        kan niet tot $2$ worden verlaagd); (b) de vier
        verzamelingen $S_i = \{x_1, \dots, x_4\} \setminus
        \{x_i\}$, voor vier punten in algemene ligging, voldoen
        aan de hypothese over drievoudige doorsneden maar niet aan
        het besluit (convexiteit telt); (c) de gesloten
        halfvlakken $H_k = \intco{k}{+\infty} \times \R$, $k \in
        \N$, snijden elkaar paarsgewijs en drievoudig, maar
        $\bigcap_k H_k = \emptyset$ (oneindige families hebben
        \omterm{def:b2:metric:compact}{compactheid} nodig).
  \item (\omterm{def:b2:metric:compact}{Compacte} Helly) Zij $(K_i)_{i \in I}$ een willekeurige
        familie \omterm{def:b2:metric:compact}{compacte} \omterm{def:b2:affine:convex}{convexe} deelverzamelingen van $\R^2$
        waarvan elke drie elkaar snijden. Toon met vraag 12 en de
        eigenschap van Borel--Lebesgue
        (\cref{thm:b2:metric:borellebesgue}) aan dat $\bigcap_{i
        \in I} K_i \neq \emptyset$.
  \item (Eerste dividend) Zij $S$ een eindige verzameling punten
        van het vlak en $r > 0$. Toon aan: liggen elke drie punten
        van $S$ in een zekere gesloten schijf met straal $r$, dan
        ligt $S$ in één gesloten schijf met straal $r$. \emph{(Pas
        Helly toe op de schijven $\overline D(p, r)$, $p \in S$.)}
\end{enumerate}

\medskip
\noindent\textbf{Deel IV --- De stelling van het centrumpunt.} Een
\emph{centrumpunt} van een eindige verzameling $S$ van $n$ punten
van het vlak is een punt $c$ (niet noodzakelijk in $S$) zodanig dat
elk gesloten halfvlak dat $c$ bevat, minstens $n/3$ punten van $S$
bevat.
\begin{enumerate}[resume]
  \item (Dimensie $1$) Toon voor reële getallen $x_1 \leq \dots
        \leq x_n$ aan dat de mediaan $c = x_{\lceil n/2 \rceil}$
        hieraan voldoet: elke gesloten halfrechte die $c$ bevat,
        bevat minstens $n/2$ van de $x_i$.
  \item (Telhulpstelling) Zijn $A, B, C$ deelverzamelingen van $S$
        met $\abs A, \abs B, \abs C > \tfrac{2n}3$, toon dan aan
        dat $A \cap B \cap C \neq \emptyset$.
  \item Zij $m = \floor{2n/3} + 1$ en zij $\mathcal F$ de
        (eindige) familie van de \omterm{def:b2:affine:convex}{convexe omhulsels}
        $\operatorname{conv}(T)$, $T \subseteq S$, $\abs T = m$.
        Toon aan dat elke drie leden van $\mathcal F$ een
        gemeenschappelijk punt hebben, en leid uit Helly een punt
        $c$ af dat aan alle gemeenschappelijk is.
  \item Bewijs dat deze $c$ een centrumpunt van $S$ is: de
        \emph{stelling van het centrumpunt}. \emph{(Bevatte een
        gesloten halfvlak door $c$ minder dan $n/3$ punten, dan zou
        zijn \omterm{def:b2:metric:topology}{open} complement een verzameling $T$ van $m$ punten
        bevatten, en zou $\operatorname{conv}(T)$ het punt $c$
        mijden.)}
  \item Scherpte: zij $n = 3k$ en plaats $k$ punten in elk van
        drie schijven met kleine straal $\varepsilon$ rond de
        hoekpunten van een grote driehoek. Toon aan dat er voor
        \emph{elk} punt $c$ van het vlak een gesloten halfvlak dat
        $c$ bevat, hoogstens $n/3$ punten van $S$ bevat, zodat de
        constante $1/3$ niet verbeterd kan worden. \emph{(Van de
        drie richtingen van $c$ naar de middelpunten van de
        schijven maken er twee een hoek van hoogstens
        $2\pi/3$.)}
\end{enumerate}

\medskip
\noindent\textbf{Deel V --- De stelling van Jung en synthese.}
\begin{enumerate}[resume]
  \item (Hulpstelling over driehoeken) Zijn $P, Q, R$ drie punten
        met paarsgewijze afstanden $\leq 1$. Toon aan dat zij in
        een gesloten schijf met straal $1/\sqrt3$ liggen.
        \emph{(Is een hoek $\geq \pi/2$, neem dan de schijf met de
        langste zijde als middellijn, met de formule voor de
        zwaartelijn $\norm{RM}^2 = \tfrac12\norm{RP}^2 +
        \tfrac12\norm{RQ}^2 - \tfrac14\norm{PQ}^2$; is de driehoek
        scherphoekig, begrens dan de omgeschreven straal $a/(2\sin
        \widehat A)$ met haar grootste hoek, die in
        $\intco{\pi/3}{\pi/2}$ ligt.)}
  \item (Jung) Leid af: elke \omterm{def:b2:metric:compact}{compacte} deelverzameling van het vlak
        met diameter $\leq 1$ is bevat in een gesloten schijf met
        straal $1/\sqrt3$.
  \item Scherpte: bewijs voor de gelijkzijdige driehoek
        $A_1A_2A_3$ met zijde $1$ en zwaartepunt $G$ de identiteit
        van Leibniz $\sum_i \norm{\vect{OA_i}}^2 =
        3\norm{\vect{OG}}^2 + \sum_i \norm{\vect{GA_i}}^2$ voor
        elk punt $O$, en besluit dat elke schijf die de drie
        hoekpunten bevat, straal $\geq 1/\sqrt3$ heeft, met
        gelijkheid alleen voor de omgeschreven schijf.
  \item (Helly in $\R^n$) Formuleer en bewijs de stelling van
        Helly in $\R^n$: zijn eindig veel \omterm{def:b2:affine:convex}{convexe verzamelingen}
        zodanig dat elke $n + 1$ ervan elkaar snijden, dan snijden
        zij elkaar alle. \emph{(Het lemma van Radon
        \cref{exo:b2:affine:7} behandelt $n + 2$ verzamelingen;
        pas dan inductie toe als in vraag 12.)}
  \item Synthese. Zet de ketting
        \[
        \text{affiene afhankelijkheid} \Rightarrow \text{Radon}
        \Rightarrow \text{Helly} \Rightarrow
        \text{centrumpunt en Jung}
        \]
        in elkaar en geef in telkens één zin aan: waar de lineaire
        algebra binnenkomt, waar de tekens van de gewichten
        binnenkomen, waar de convexiteit binnenkomt, en welke ene
        stap de dimensie van het vlak gebruikte. Wat worden de
        constanten $3$ (in Helly), $1/3$ (centrumpunt) en
        $1/\sqrt3$ (Jung) in $\R^n$? (Formuleer zonder bewijs.)
\end{enumerate}
\end{problem}
